\section{What \texorpdfstring{$\mathrm{OutTrans}{=}1$}{OutTrans=1} Is, and Why One Tensor Still Resists}\label{sec:outtrans}

The reconstruction reported in \S\ref{sec:reconstruct} left seven weight roles per layer, of which
five decoded correctly and two did not. This section resolves one of the two, identifies what the
$\mathrm{OutTrans}$ flag physically does, narrows the remaining failure to a single tensor, and
records the instruments that had to be discarded along the way.

\subsection{A validated instrument}\label{sec:oracle}

Every earlier attempt to score a candidate weight ordering used a statistic that was \emph{assumed}
to work. One of them was later shown to be invalid outright (\S\ref{sec:retract-basis}). The
instrument used here is derived from the computation rather than from weight statistics, and is
calibrated against a model whose ordering is known to be correct before being applied to the
recovered weights.

The attention output projection $O$ writes residual position $j$; the FFN gate reads it. The composed
map therefore contracts the residual index,
\[
  \mathrm{residual}_j=\textstyle\sum_i \mathrm{attn}_i\,O[j,i],\qquad
  \mathrm{ffn}_k=\textstyle\sum_j \mathrm{residual}_j\,\mathrm{gate}[k,j],
\]
so a mis-ordered residual axis degrades $\lVert \mathrm{gate}\cdot O\rVert_F$. Scoring a candidate
permutation $\pi$ against random re-pairings gives a $z$-score. (This is a quadratic-assignment
objective with \emph{raw} Gram matrices; earlier attempts used row-normalised cosine matrices, which
discard exactly the magnitude information the objective depends on---a likely reason they found
nothing.)

\begin{finding}[the instrument is calibrated, not assumed]\label{find:oracle}
On a correctly-ordered reference model of the same architecture family (RMSNorm, per-head QK-norm,
GQA, SwiGLU), the true pairing separates from random re-pairings at $z\approx+300$ to $+400$
depending on the pairing type. Only then was it applied to the recovered weights.
\end{finding}
\noindent\emph{Evidence.} Reference control, averaged over layers: $O\!\to\!\mathrm{gate}$ (direct)
$+406$; $\mathrm{up}\!\to\!\mathrm{down}$ (SwiGLU between) $+407$;
$\mathrm{down}\!\to\!\mathrm{gate}(L{+}1)$ (an attention block between) $+328$;
$\mathrm{gate}\!\to\!\mathrm{down}$ $+57$. Every pairing type is informative, including across
nonlinearities.

\subsection{The attention output projection is stored transposed}

\begin{finding}[$O$ is stored transposed]\label{find:otrans}
The attention output projection is stored with its axes swapped relative to the
$\mathrm{OutTrans}{=}0$ roles: the axis the decoder treats as the input is in fact the residual
(output) axis.
\end{finding}
\noindent\emph{Evidence.} Three mutually independent methods agree.
(i) \emph{Composition}: as decoded, $z=+2.9$ (random); transposed, $z=+96$ to $+105$.
(ii) \emph{Attention-head signature}: one axis of $O$ must be the head-concatenation
($16\times64$), which shows $64$-periodic block structure. Calibrating on tensors whose orientation
is not in doubt gives a head axis at $+0.0217$ and a residual axis at $-0.0005$; $O$'s two axes
measure $+0.0195$ and $-0.0015$ respectively---i.e.\ reversed.
(iii) \emph{Per-role census}: scoring every role against a verified residual reference,
$Q$ $+300.7$, $K$ $+184.9$, $V$ $+174.5$, $\mathrm{up}$ $+523.9$ all confirm the convention
(input $=$ residual); $O$ alone violates it, at $+160.1$ on the opposite axis.

This reduces the open problem from two roles to one: with $O$ transposed, six of seven roles per
layer are positively confirmed.

\subsection{What the flag physically does}

The ANE task descriptors carry L2 buffer addressing---base addresses and per-axis strides---which
describes the \emph{activation} layout directly.

\begin{finding}[$\mathrm{OutTrans}{=}1$ is a strided, interleaved output write]\label{find:outtrans2}
The flag is perfectly determined by the L2 result stride: $\mathrm{OutTrans}{=}0$ writes channels
contiguously, $\mathrm{OutTrans}{=}1$ writes them $16\times$ strided, with the tasks of one tensor
interleaved into each block.
\end{finding}
\noindent\emph{Evidence.} Over all $988$ weight-bearing tasks in the shipped binary: every one of the
$180$ $\mathrm{OutTrans}{=}1$ tasks has a large result stride (\code{0x810} or \code{0x800}); none of
the $808$ $\mathrm{OutTrans}{=}0$ tasks does---those use \code{0x80}, exactly one $64$-wide
\textrm{fp16} row. For the down-projection specifically, the four tasks read their FFN inputs
contiguously (\code{SrcStrideC}${=}$\code{0x80}) but write to bases \code{0x14c040}, \code{0x14c240},
\code{0x14c440}, \code{0x14c640}---\code{0x200} apart, i.e.\ four channels---under a \code{0x810}
per-channel stride.

Two consequences follow. First, the flag governs activation storage, not coefficient encoding, which
is consistent with the positional read below. Second, the down-projection's \emph{input} axis is not
permuted at the activation level; only its output side is in question.

\subsection{The coefficient layout, read in the shipped configuration}

The shipped tiles use a $128$-byte coefficient header (\code{CoeffSize} \code{0x6480}) where a
straightforwardly compiled convolution emits $64$ bytes (\code{0x6440}). That mode initially could
not be produced: the per-channel-scale palettization that seemed to explain it lowers to an operator
the ANE compiler rejects.

\begin{finding}[the extra header is a bias, and the payload layout is unchanged]\label{find:bias}
The $128$-byte header is emitted when the convolution carries a \emph{bias}, not a per-channel scale.
With that mode forced, a positional read reproduces the identical coefficient map as in the
$64$-byte mode.
\end{finding}
\noindent\emph{Evidence.} Compiling with $\mathrm{bias}[o]=100+o$ and reading the header back gives
\code{100 101 \dots 115} for bank~$0$ and \code{116 \dots 131} for bank~$1$---an exact match. A
positional read at the shipped geometry ($3200\to256$, $16$ banks) then recovers a \emph{perfect
bijection} over all $819{,}200$ positions,
\[
  o = 16\,b + (\mathrm{slot} \bmod 16), \qquad i = \lfloor \mathrm{slot}/16 \rfloor,
\]
identical in both header modes. The shipped tiles' geometry, codec and slot decomposition are
therefore confirmed against hardware ground truth, not assumed.

(The shipped values in that header slot are all positive and tightly banded, $0.07$--$0.44$, against
a codebook spanning $\pm0.8$---the signature of a per-output scale rather than a bias. The slot is
evidently reused by mode, and the decoder's interpretation is retained.)

\subsection{The remaining defect, and why it is hard}

\begin{finding}[the failure is confined to one tensor]\label{find:downonly}
With $O$ transposed, every weight pairing that does not involve the FFN down-projection is aligned,
and every pairing that does involve it sits at the random level.
\end{finding}
\noindent\emph{Evidence.} $O\!\to\!\mathrm{gate}$ $+104.7$ and $Q\!\to\!O$ $+70.7$ (aligned) against
$\mathrm{down}\!\to\!Q(L{+}1)$ $+1.4$, $\mathrm{down}\!\to\!\mathrm{gate}(L{+}1)$ $+1.3$,
$\mathrm{gate}\!\to\!\mathrm{down}$ $-2.6$, $\mathrm{up}\!\to\!\mathrm{down}$ $-1.4$. The two tests
are independent: permuting the tensor's columns cannot change $\lVert\mathrm{gate}\cdot
\mathrm{down}\rVert_F$, and permuting its rows cannot change the other, so both axes are implicated
separately.

The tensor's \emph{values} are correct. Permutation preserves singular values, and the recovered
down-projection has a clearly trained spectrum ($s_1/s_{10}=2.11$, $s_1/s_{100}=2.90$) against a
random control at $1.02$ and $1.14$. The defect is ordering alone.

Reading the channel assignment out of the L2 addresses gives the right direction but not the answer.
The thirteen tasks building one FFN tensor write to \emph{descending} addresses as the task index
rises, spanning \code{0x064000}--\code{0x0c8000} $=3200\times128$\,B exactly, so the file order is
\emph{reversed} relative to channel order and the $128$-wide half-block---first in the file---holds
the \emph{last} $128$ channels. Applying that map improves
$\mathrm{gate}\!\to\!\mathrm{down}$ from $+0.3$ to $+2.3$, and the analogous interleave improves
$\mathrm{down}\!\to\!\mathrm{gate}(L{+}1)$ from $+0.4$ to $+3.0$. Both improve consistently across
layers, and both remain an order of magnitude short of the calibrated reference. The likely obstacle
is arithmetic the traces do not pin down: the stride \code{0x810} is $16\times$\code{0x80}${}+{}$\code{0x10},
so channel slots are not uniformly spaced and any clean closed form is an approximation.

\begin{observation}[the constraint system is underdetermined]\label{obs:underdet}
The $3200$-wide FFN channel axis is referenced by only four tensors---$\mathrm{gate}$ and
$\mathrm{up}$ (outputs), $\mathrm{down}$ (input), and the corresponding adapter factor---all of which
are among the tensors in question. The residual axis, by contrast, is touched by seven roles, which
is why it could be pinned. No external anchor for the FFN ordering exists in the shipped data.
\end{observation}

\subsection{The compiler is not the variable}

Because the ANE compile is performed by \code{ANECompiler.framework}, which ships with the operating
system rather than with the developer toolchain, it was natural to expect a different OS build to
emit the configuration needed to read the layout directly. It does not.

\begin{finding}[compiler version and target do not change the outcome]\label{find:compiler}
Across two macOS builds spanning a major \code{ANECompiler} version difference, two toolchain
versions, six ANE architecture targets and two optimisation settings, no configuration emits
$\mathrm{OutTrans}{=}1$ on a convolution that carries coefficients.
\end{finding}
\noindent\emph{Evidence.} \code{ANECompiler} $10.24.3$ (macOS~$27.0$) and $9.509.0$
(macOS~$26.5.2$) produce byte-identical task structure on all five probe graphs: the flag appears on
$1$, $2$ and $4$ tasks for the transpose and attention graphs and is weight-bearing on \emph{none}.
Two toolchain versions on one host likewise agree exactly, including identical failures---expected,
since the toolchain supplies only the front end. Older ANE targets (h14, h15) never emit the flag at
all; h16--h18 emit it only on weightless shuffle tasks. Simulator runtimes cannot substitute: the
simulator has no Neural Engine, ships no ANE frameworks, and its own weight files are suffixed
\code{\_nonane}.

The difference in the shipped binary---which contains $180$ weight-bearing $\mathrm{OutTrans}{=}1$
tasks---therefore originates upstream of the compiler, in how the graph reaches it.

\subsection{Instruments discarded}\label{sec:retract-basis}

Two more instruments were withdrawn during this work, in both cases because a null result was being
read as evidence when the statistic had never been shown to carry information.

\textbf{The residual-magnitude oracle.} A statistic scoring whether a writer's per-channel magnitude
tracks a reader's was used to localise the defect. Calibrated afterwards on the correctly-ordered
reference model, the reader-to-reader baseline it depends on measures $-0.176$---writer and reader
magnitudes simply do not track each other in a correct model, so the recovered weights scoring $\approx 0$
was the expected result rather than a defect. Every localisation drawn from it was withdrawn.

\textbf{Adapter data as ground truth.} The shipped adapter constant-data file is loaded verbatim into
an unbacked mutable kernel segment (its size matches the segment's VM size exactly, and the file is
layer-major with autocorrelation $0.893$ at that period), which made it attractive as
quantization-free ground truth spanning both $\mathrm{OutTrans}$ modes. Its internal role offsets
were re-derived exactly---segmenting by RMS \emph{level} rather than by jump, then requiring the run
lengths to be a permutation of the architecturally fixed multiset, which recovers $14$ segments and
corrects the role order---but the file still cannot be validated against the base weights, because no
controllable test exists for whether a randomly-initialised adapter factor \emph{should} align with
the weight it adapts. Conclusions drawn from it are therefore not used.

\subsection{Status}

Six of seven weight roles per layer are positively confirmed rather than assumed. The
down-projection's values are correct, its geometry, codec and slot decomposition are verified against
hardware ground truth in the shipped configuration, and its channel layout is partially recovered
from the shipped addressing. A coherent forward pass is still not achieved: one mis-ordered
residual-writing projection is sufficient to destroy it, and closing the gap requires either the ANE
tiling specification or a weight-bearing $\mathrm{OutTrans}{=}1$ compile, neither of which is
reachable from the shipped assets and available toolchains.
